% ============================================================
%  Hands-On 4 -- Implementasi Multi-threading
%  Praktikum Sistem Operasi
% ============================================================

\chapter{Implementasi Multi-threading}
\label{chap:handson-04}

% ----------------------------------------------------------------
\section{Tujuan}

Jawablah pertanyaan-pertanyaan berikut berdasarkan hasil praktikum Anda
sesuai urutan sesi pada Modul~8.
Sertakan bukti berupa \textit{screenshot} atau penggalan keluaran
terminal yang relevan pada setiap jawaban.

% ----------------------------------------------------------------
\section{Sesi 1 --- Persiapan Lingkungan Kerja}

\subsection{Masuk ke Container Docker}

Jalankan perintah berikut di terminal, kemudian masuk ke direktori
\textit{workspace}:

\begin{lstlisting}[language=bash, caption={Menjalankan container dan masuk workspace}]
docker run --rm -it \
    -v "C:\path\ke\starter-code:/workspace" \
    sinanonym/ifos-env
cd /workspace
\end{lstlisting}

\begin{enumerate}

    \item \textbf{Compile starter code sebelum implementasi.}\\
    Jalankan perintah berikut di dalam container:

    \begin{lstlisting}[language=bash, caption={Perintah kompilasi awal}]
make clean && make
    \end{lstlisting}

    Tuliskan seluruh \textit{warning} yang muncul pada saat kompilasi dan
    jelaskan mengapa \textit{warning} tersebut muncul pada kode
    \textit{stub}.\\[0.3cm]
    \textbf{Jawaban:}\\
    \textit{[Isi jawaban di sini]}

    \item \textbf{Jalankan QEMU sebelum implementasi.}\\
    Jalankan perintah berikut lalu sertakan \textit{screenshot} layar QEMU.
    Jelaskan mengapa sistem berhenti di sana dan belum menampilkan
    output Thread~A/B/C.

    \begin{lstlisting}[language=bash, caption={Menjalankan emulator QEMU}]
make run
    \end{lstlisting}

    \textbf{Screenshot:}
    \begin{figure}[H]
        \centering
        \fbox{\begin{minipage}{0.92\linewidth}
            \centering\vspace{2.5cm}
            \textit{[Tempel screenshot layar QEMU di sini]}
            \vspace{2.5cm}
        \end{minipage}}
        \caption{Tampilan QEMU sebelum implementasi}
    \end{figure}

    \textbf{Jawaban:}\\
    \textit{[Isi jawaban di sini]}

\end{enumerate}

% ----------------------------------------------------------------
\section{Sesi 2 --- Implementasi \texttt{thread\_create()}}

\subsection{Thread Control Block dan Layout Memori}

\begin{enumerate}

    \item \textbf{Field-field \texttt{struct thread}.}\\
    Perhatikan definisi \texttt{struct thread} di \texttt{src/thread.h}.
    Isilah tabel berikut berdasarkan kode sumber.

    \begin{table}[!ht]
    \caption{Field-field pada \texttt{struct thread}}
    \label{tab:h04-struct-thread}
    \centering
    \begin{tabulary}{\linewidth}{|L|L|L|}
    \hline
    \textbf{Field} & \textbf{Tipe Data} & \textbf{Fungsi} \\
    \hline
    \texttt{tid}      & \texttt{tid\_t}              & \textit{[Isi]} \\
    \hline
    \texttt{status}   & \texttt{enum thread\_status} & \textit{[Isi]} \\
    \hline
    \texttt{name}     & \texttt{char[16]}            & \textit{[Isi]} \\
    \hline
    \texttt{stack}    & \texttt{uint8\_t *}          & \textit{[Isi]} \\
    \hline
    \texttt{priority} & \texttt{int}                 & \textit{[Isi]} \\
    \hline
    \texttt{magic}    & \texttt{unsigned}            & \textit{[Isi]} \\
    \hline
    \end{tabulary}
    \end{table}

    \item \textbf{Layout halaman 4~KB setiap thread.}\\
    Setiap thread menempati satu halaman (\textit{page}) memori berukuran
    4~KB. Isilah kolom \textit{Isi} pada tabel layout di bawah ini
    untuk menggambarkan posisi \texttt{struct thread}, \textit{stack},
    dan ketiga frame yang disiapkan oleh \texttt{thread\_create()}.

    \begin{table}[!ht]
    \caption{Layout halaman 4~KB milik sebuah thread}
    \label{tab:h04-page-layout}
    \centering
    \begin{tabulary}{\linewidth}{|C|L|}
    \hline
    \textbf{Offset (relatif)} & \textbf{Isi} \\
    \hline
    \texttt{+4096} & \textit{[Isi: puncak halaman / ESP awal]} \\
    \hline
    \texttt{...}   & \texttt{switch\_threads\_frame} \\
    \hline
    \texttt{...}   & \texttt{switch\_entry\_frame} \\
    \hline
    \texttt{...}   & \texttt{kernel\_thread\_frame} \\
    \hline
    \texttt{...}   & \textit{[Isi: area stack kosong, tumbuh ke bawah ($\downarrow$)]} \\
    \hline
    \texttt{+0}    & \texttt{struct thread} (TCB) \\
    \hline
    \end{tabulary}
    \end{table}

    \item \textbf{Analisis fungsi \texttt{running\_thread()}.}\\
    Perhatikan potongan kode berikut:

    \begin{lstlisting}[language=c, caption={Implementasi \texttt{running\_thread()} di \texttt{src/thread.c}}]
struct thread *running_thread(void) {
    uint32_t esp;
    __asm__ volatile ("mov %%esp, %0" : "=g" (esp));
    return (struct thread *)(void *)(esp & ~(uint32_t)(PGSIZE - 1));
}
    \end{lstlisting}

    \begin{enumerate}
        \item[(a)] Jelaskan apa yang dilakukan operasi bitwise
                   \texttt{esp \& \textasciitilde(PGSIZE - 1)} dan
                   mengapa hasilnya selalu merupakan alamat awal
                   \textit{page}.\\[0.3cm]
                   \textbf{Jawaban:}\\
                   \textit{[Isi jawaban di sini]}

        \item[(b)] Mengapa teknik ini hanya benar jika
                   \texttt{struct thread} selalu berada di offset~0
                   dari halamannya? Apa dampaknya jika
                   \texttt{initial\_thread\_page} di \texttt{boot.asm}
                   tidak di-\textit{align} ke 4096~byte?\\[0.3cm]
                   \textbf{Jawaban:}\\
                   \textit{[Isi jawaban di sini]}

        \item[(c)] Apa fungsi \texttt{THREAD\_MAGIC} dan bagaimana
                   \texttt{thread\_current()} menggunakannya untuk
                   mendeteksi \textit{stack overflow}?\\[0.3cm]
                   \textbf{Jawaban:}\\
                   \textit{[Isi jawaban di sini]}
    \end{enumerate}

\end{enumerate}

\subsection{Implementasi Fungsi \texttt{thread\_create()}}

\begin{enumerate}

    \item \textbf{Tiga stack frame pada \texttt{thread\_create()}.}\\
    Isilah tabel berikut berdasarkan komentar petunjuk di
    \texttt{src/thread.c}.

    \begin{table}[!ht]
    \caption{Tiga stack frame yang disiapkan \texttt{thread\_create()}}
    \label{tab:h04-stack-frames}
    \centering
    \begin{tabulary}{\linewidth}{|C|L|L|L|}
    \hline
    \textbf{Urutan push} & \textbf{Nama frame} &
    \textbf{Field penting} & \textbf{Nilai yang diisi} \\
    \hline
    1 --- paling atas & \texttt{kernel\_\allowbreak thread\_\allowbreak frame} &
        \texttt{eip}, \texttt{function}, \texttt{aux} &
        \textit{[Isi]} \\
    \hline
    2 --- tengah & \texttt{switch\_\allowbreak entry\_\allowbreak frame} &
        \texttt{eip} &
        \textit{[Isi]} \\
    \hline
    3 --- paling bawah & \texttt{switch\_\allowbreak threads\_\allowbreak frame} &
        \texttt{eip}, register \textit{callee-saved} &
        \textit{[Isi]} \\
    \hline
    \end{tabulary}
    \end{table}

    \item \textbf{Alur eksekusi pertama thread baru.}\\
    Saat thread baru \textbf{pertama kali} dijadwalkan melalui
    \texttt{switch\_threads()}, jelaskan secara berurutan apa yang
    terjadi mulai dari \texttt{switch\_entry} hingga fungsi thread
    (\texttt{func\_a}, \texttt{func\_b}, atau \texttt{func\_c})
    benar-benar mulai berjalan.\\[0.3cm]
    \textbf{Jawaban:}\\
    \textit{[Isi jawaban di sini]}

    \item \textbf{Tulis implementasi \texttt{thread\_create()}.}\\
    Salin kode yang telah Anda tulis di \texttt{src/thread.c}:

    \begin{lstlisting}[language=c, caption={Implementasi \texttt{thread\_create()} oleh mahasiswa}]
tid_t thread_create(const char *name, int priority,
                    thread_func *function, void *aux) {

    /* Tulis implementasi Anda di sini */

}
    \end{lstlisting}

    \item \textbf{Penanganan kegagalan alokasi.}\\
    Kapan \texttt{allocate\_thread\_page()} dapat mengembalikan
    \texttt{NULL}, dan bagaimana \texttt{thread\_create()} harus
    menanganinya?\\[0.3cm]
    \textbf{Jawaban:}\\
    \textit{[Isi jawaban di sini]}

\end{enumerate}

% ----------------------------------------------------------------
\section{Sesi 3 --- Implementasi \textit{Cooperative Round-Robin}}

\subsection{Konsep Status Thread}

\begin{enumerate}

    \item \textbf{Perbedaan \texttt{THREAD\_READY} dan
    \texttt{THREAD\_BLOCKED}.}\\
    Lengkapi tabel perbandingan berikut.

    \begin{table}[!ht]
    \caption{Perbandingan status \texttt{THREAD\_READY} dan \texttt{THREAD\_BLOCKED}}
    \label{tab:h04-thread-status}
    \centering
    \begin{tabulary}{\linewidth}{|L|L|L|}
    \hline
    \textbf{Aspek} & \textbf{THREAD\_READY} & \textbf{THREAD\_BLOCKED} \\
    \hline
    Siapa yang menetapkan? & \textit{[Isi]} & \textit{[Isi]} \\
    \hline
    Masuk ke \textit{ready queue}? & \textit{[Isi]} & \textit{[Isi]} \\
    \hline
    Cara kembali ke \texttt{RUNNING} & \textit{[Isi]} & \textit{[Isi]} \\
    \hline
    Contoh penggunaan & \textit{[Isi]} & \textit{[Isi]} \\
    \hline
    \end{tabulary}
    \end{table}

\end{enumerate}

\subsection{Implementasi Fungsi \textit{Scheduling}}

\begin{enumerate}

    \item \textbf{Tulis implementasi \texttt{thread\_yield()}.}\\
    Salin kode yang telah Anda tulis di \texttt{src/thread.c}:

    \begin{lstlisting}[language=c, caption={Implementasi \texttt{thread\_yield()} oleh mahasiswa}]
void thread_yield(void) {

    /* Tulis implementasi Anda di sini */

}
    \end{lstlisting}

    \item \textbf{Tulis implementasi \texttt{next\_thread\_to\_run()}
    (Round-Robin FIFO).}\\
    Salin kode yang telah Anda tulis di \texttt{src/thread.c}:

    \begin{lstlisting}[language=c, caption={Implementasi \texttt{next\_thread\_to\_run()} oleh mahasiswa}]
struct thread *next_thread_to_run(void) {

    /* Tulis implementasi Anda di sini */

}
    \end{lstlisting}

    \item \textbf{Mengapa mengembalikan \texttt{initial\_thread},
    bukan \texttt{NULL}?}\\
    Jelaskan apa yang akan terjadi pada fungsi \texttt{schedule()} jika
    \texttt{next\_thread\_to\_run()} mengembalikan \texttt{NULL} saat
    antrian kosong.\\[0.3cm]
    \textbf{Jawaban:}\\
    \textit{[Isi jawaban di sini]}

    \item \textbf{\texttt{thread\_block()} vs.\ \texttt{thread\_yield()}
    di akhir \texttt{kernel\_main}.}\\
    Jelaskan perbedaan perilaku kedua fungsi tersebut setelah Thread~A,
    B, dan C berhasil dibuat, dan mengapa \texttt{thread\_block()}
    adalah pilihan yang benar.\\[0.3cm]
    \textbf{Jawaban:}\\
    \textit{[Isi jawaban di sini]}

\end{enumerate}

% ----------------------------------------------------------------
\section{Sesi 4 --- Kompilasi dan Pengujian QEMU}

\begin{enumerate}

    \item \textbf{Compile ulang setelah implementasi selesai.}

    \begin{lstlisting}[language=bash, caption={Perintah build dan run setelah implementasi}]
make clean && make && make run
    \end{lstlisting}

    Sertakan \textit{screenshot} layar QEMU yang menampilkan ketiga
    thread berjalan bergantian.

    \begin{figure}[H]
        \centering
        \fbox{\begin{minipage}{0.92\linewidth}
            \centering\vspace{2.5cm}
            \textit{[Tempel screenshot layar QEMU setelah implementasi di sini]}
            \vspace{2.5cm}
        \end{minipage}}
        \caption{Tampilan QEMU setelah implementasi berhasil}
    \end{figure}

    \item \textbf{Verifikasi urutan penjadwalan.}\\
    Ekspektasi keluaran sistem yang benar adalah sebagai berikut:

    \begin{lstlisting}[caption={Ekspektasi output Round-Robin di layar QEMU}]
===================================
  IF-OS  |  Praktikum Sistem Operasi
===================================

Inisialisasi sistem thread...
Membuat thread A, B, C...
Memulai scheduling...

[Thread A] tick 0
[Thread B] tick 0
[Thread C] tick 0
[Thread A] tick 1
[Thread B] tick 1
[Thread C] tick 1
...
    \end{lstlisting}

    Apakah urutan Thread~A $\to$ Thread~B $\to$ Thread~C $\to$
    Thread~A $\to \ldots$ terjaga secara konsisten pada hasil Anda?
    Jelaskan mengapa urutan ini dapat dijamin oleh implementasi FIFO
    pada \textit{ready queue}.\\[0.3cm]
    \textbf{Jawaban:}\\
    \textit{[Isi jawaban di sini]}

    \item \textbf{Analisis \textit{warning} compiler sebelum dan setelah
    implementasi.}\\
    Apakah \textit{warning} \texttt{unused-function} dan
    \texttt{unused-variable} yang muncul sebelum implementasi masih
    ada setelah Anda mengimplementasikan ketiga fungsi? Jelaskan
    mengapa \textit{warning} tersebut hilang.\\[0.3cm]
    \textbf{Jawaban:}\\
    \textit{[Isi jawaban di sini]}

\end{enumerate}

% ----------------------------------------------------------------
\section{Eksplorasi Mandiri}

\begin{enumerate}

    \item \textbf{Menambah Thread keempat.}\\
    Tambahkan Thread~D dengan fungsi dan warna VGA berbeda di
    \texttt{src/kernel.c}, kemudian compile ulang dan jalankan.

    \begin{enumerate}
        \item[(a)] Sertakan \textit{screenshot} layar QEMU yang
                   menampilkan empat thread berjalan bergantian.

        \begin{figure}[H]
            \centering
            \fbox{\begin{minipage}{0.92\linewidth}
                \centering\vspace{2.5cm}
                \textit{[Tempel screenshot empat thread di sini]}
                \vspace{2.5cm}
            \end{minipage}}
            \caption{Tampilan QEMU dengan empat thread}
        \end{figure}

        \item[(b)] Berapa nilai \texttt{MAX\_THREADS} saat ini?
                   Apa yang terjadi jika Anda membuat thread melebihi
                   batas tersebut? Tunjukkan dengan percobaan dan
                   sertakan bukti.\\[0.3cm]
                   \textbf{Jawaban:}\\
                   \textit{[Isi jawaban di sini]}
    \end{enumerate}

    \item \textbf{Simulasi \textit{finite thread}.}\\
    Ubah \texttt{func\_a} di \texttt{src/kernel.c} agar hanya
    mencetak \textbf{5 kali} lalu berhenti (ganti \texttt{while(1)}
    dengan loop terbatas). Amati apakah Thread~B dan Thread~C tetap
    berjalan setelah Thread~A selesai, lalu jelaskan peran
    \texttt{thread\_exit()} dalam skenario ini.\\[0.3cm]
    \textbf{Jawaban:}\\
    \textit{[Isi jawaban di sini]}

    \item \textbf{Refleksi: \textit{cooperative} vs.\
    \textit{preemptive scheduling}.}\\
    Sistem IF-OS saat ini menggunakan \textit{cooperative scheduling}
    di mana thread hanya berganti saat secara sadar memanggil
    \texttt{thread\_yield()}.

    \begin{enumerate}
        \item[(a)] Apa kelemahan utama \textit{cooperative scheduling}
                   dibandingkan \textit{preemptive scheduling}?
                   Berikan contoh skenario nyata di mana kelemahan
                   ini berbahaya.\\[0.3cm]
                   \textbf{Jawaban:}\\
                   \textit{[Isi jawaban di sini]}

        \item[(b)] Komponen perangkat keras apa yang perlu ditambahkan
                   agar kernel dapat memaksa pergantian thread secara
                   otomatis tanpa bergantung pada kepatuhan thread
                   itu sendiri?\\[0.3cm]
                   \textbf{Jawaban:}\\
                   \textit{[Isi jawaban di sini]}
    \end{enumerate}

\end{enumerate}
